% Centri — plan from the teaching review of the tiered material (2026-07-23)
% Standalone deck: every slide readable cold, no narration needed.
% Compile: pdflatex -interaction=nonstopmode centri-plan-2026-07-23.tex  (run twice)
\documentclass[aspectratio=169,10pt]{beamer}
\usetheme{metropolis}
\metroset{sectionpage=progressbar, subsectionpage=none, numbering=fraction, progressbar=frametitle}

\usepackage{booktabs}
\usepackage{amsmath,amssymb}
\usepackage{xcolor}
\usepackage{colortbl}
\usepackage{graphicx}
\graphicspath{{figs/}}
\usepackage[most]{tcolorbox}

% ---- palette (shared with the other Centri decks) ------------------------
\definecolor{cMeas}{HTML}{1F6FB2}
\definecolor{cPed}{HTML}{B5651D}
\definecolor{cApp}{HTML}{2E7D32}
\definecolor{cInfra}{HTML}{5E35B1}
\definecolor{cAccent}{HTML}{C62828}
\definecolor{cGrey}{HTML}{616161}
\definecolor{cBasic}{HTML}{2E7D32}
\definecolor{cInter}{HTML}{1F6FB2}
\definecolor{cAdv}{HTML}{5E35B1}
\setbeamercolor{frametitle}{bg=cPed}
\setbeamercolor{progress bar}{fg=cAccent}

\newcommand{\ok}{{\color{cApp}\textbf{\checkmark}}}
\newcommand{\nope}{{\color{cGrey}--}}
\newcommand{\nw}{{\scriptsize\color{cAccent}\textbf{new}}}
\newcommand{\takeaway}[1]{\vfill{\footnotesize\itshape\color{cGrey}#1\par}}

\title{Centri --- Reworking the Three Levels}
\subtitle{What a physics-teaching review of the generated material changes, and in what order}
\author{Damar Albaribin}
\institute{MSc Thesis --- advisor: Prof.\ Wu-Yuin Hwang}
\date{2026-07-23 \quad {\scriptsize\color{cGrey}rough plan for discussion}}

\begin{document}
\maketitle

% =========================================================================
\begin{frame}{Where this comes from}
\small
\textbf{Centri} turns a phone video of something spinning into physics worksheets at \textbf{three
levels} --- basic, intermediate, advanced --- each with annotated figures and worked examples.
All three are written from the \emph{same} set of measured numbers, so they differ in demand
without disagreeing on fact.
\vspace{0.5em}

Those worksheets have now been read by someone who has taught physics in a classroom
(\textbf{Vinsa}, who also authors the sibling app \textbf{P-MAGIC}). The verdict, in short:
\vspace{0.4em}

\begin{tcolorbox}[colback=cPed!6,colframe=cPed,boxsep=2pt,left=6pt,right=6pt,top=3pt,bottom=3pt]
\small The \textbf{basic} level is fine. The \textbf{intermediate} level reaches for equations
before the reader has understood the picture. The \textbf{advanced} level is written in language a
high-school student will not have met --- calibration, scale-free quantities, Jensen's inequality.
\end{tcolorbox}
\vspace{0.5em}

\small Read alongside it: the ladder from one level to the next is \textbf{too steep}. Each level is
a single jump rather than a few graded steps, and nothing carries a reader from the end of one
level into the start of the next.
\takeaway{The architecture is right and stays; what changes is the reading level of the top rung and
the size of the steps between rungs.}
\end{frame}

% =========================================================================
\begin{frame}{A defect the review exposed}
\scriptsize
Checking one review note --- \emph{``why is the turn rate negative, and where do we say so?''} ---
turned up a live error in the two ceiling-fan clips.

\begin{center}
\includegraphics[width=0.80\textwidth]{fan-sign-4028.png}
\end{center}

\begin{columns}[T]
\begin{column}{0.32\textwidth}
A minus sign only records \textbf{which way round} something turns; these two fans turn the
direction recorded as negative.
\end{column}
\begin{column}{0.34\textwidth}
``Speeding up'' is read from the \textbf{sign of the slope} --- and a negative value climbing
towards zero slopes upward.
\end{column}
\begin{column}{0.32\textwidth}
One fan's \textbf{basic} worksheet teaches \emph{``speeding up motion''} while its own graph reads
``slowing down''.
\end{column}
\end{columns}
\takeaway{Two of seven clips affected. The fix: judge the trend from the speed, not the signed value.}
\end{frame}

% =========================================================================
\begin{frame}{Lowering the reading level}
\scriptsize
Three phrases were named as beyond a high-school reader. Each is in the student worksheet because
\textbf{our own specification asks for it} --- so each is a one-line change, not a model limitation.
\vspace{4pt}

\renewcommand{\arraystretch}{1.45}
\begin{tabular}{@{}p{4.5cm}p{4.6cm}p{4.6cm}@{}}
\toprule
\textbf{term the student meets today} & \textbf{what it is really saying} & \textbf{where it goes} \\
\midrule
``which quantities are independent of the \textbf{calibration}'' &
angle measurements do not depend on how we sized the scene; metre values do &
said in plain words to the student \\
``the angular quantities are \textbf{scale-free}'' &
the same idea, named twice &
merged into the sentence above \\
``formally $\langle\omega^2\rangle \ge \langle\omega\rangle^2$, a \textbf{Jensen inequality}'' &
squaring an average is not averaging the squares &
\textbf{teacher copy only} \\
\bottomrule
\end{tabular}
\vspace{6pt}

\small A separate \textbf{teacher copy} of every worksheet already exists and already receives the
answer key, so this is re-routing rather than new machinery.
\vspace{2pt}

Alongside it, the \textbf{intermediate} level is reordered: read the graph in words \emph{first},
then meet the formula as the compact way to say what the reader has already seen.
\takeaway{Keep the idea, drop the vocabulary, and give the formal version to the teacher.}
\end{frame}

% =========================================================================
\begin{frame}{Smaller steps, and bridges}
\scriptsize
The proposal: three graded steps inside each level, and a bridge carrying the reader across.
\vspace{4pt}

\renewcommand{\arraystretch}{1.5}
\begin{tabular}{@{}clll@{}}
\toprule
& \textbf{\color{cBasic}A --- basic} & \textbf{\color{cInter}B --- intermediate} &
\textbf{\color{cAdv}C --- advanced} \\
\midrule
\textbf{1} & which quantities are used &
{\color{cAccent}read the graph in words,} \emph{then} the equation &
compare two moments \\
\textbf{2} & what each one means & put this clip's numbers in &
compare the phases of the motion \\
\textbf{3} & which equation joins them &
{\color{cAccent}how long to slow from this speed to that one} &
{\color{cAccent}which claims does this video actually support?} \\
\midrule
\textbf{bridge} & \multicolumn{3}{l}{A3 $\rightarrow$ B1 and B3 $\rightarrow$ C1 --- one staircase,
not three separate documents} \\
\bottomrule
\end{tabular}
\vspace{2pt}

{\scriptsize\color{cAccent}Red = proposed here, not yet agreed. B1 folds in the ``interpret before
you calculate'' note; B3 is the ``squeeze the video data harder'' idea; C3 keeps the honesty
argument as a \emph{task} rather than as vocabulary.}
\vspace{6pt}

\begin{columns}[T]
\begin{column}{0.49\textwidth}
\textbf{\color{cAdv}Advanced stops being ``intermediate with more numbers.''}
Its distinguishing job becomes \emph{comparison and judgement} --- how long to slow from this speed
to that one, one phase against another, one scenario against another.
\end{column}
\begin{column}{0.49\textwidth}
\textbf{\color{cApp}The measured difficulty ladder stays the target.}
Reading grade 7.8 / 8.1 / 11.7 and ideas-held-at-once few / more / many are measured from the text,
not asserted --- the sub-steps have to keep that ladder rising smoothly.
\end{column}
\end{columns}
\takeaway{The thinking level is what separates the tiers; the sub-steps make each tier reachable
from the one below.}
\end{frame}

% =========================================================================
\begin{frame}{Showing the correction, not describing it}
\scriptsize
The worksheet currently draws \textbf{one} line and explains the correction in words. Draw
\textbf{both} --- what came off the video, and what we corrected it to. Mock-up on a real clip:

\begin{center}
\includegraphics[width=0.68\textwidth]{wheel-ac-corrected.png}
\end{center}
\vspace{-4pt}

\begin{columns}[T]
\begin{column}{0.49\textwidth}
The wheel is filmed at a slant, so the handle appears to race and stall once per turn. That is the
\textbf{camera position}, not the wheel --- and it lands in the inward pull, which follows the turn
rate \emph{squared}.
\end{column}
\begin{column}{0.49\textwidth}
Same video, both lines. A reader can \emph{see} what was removed and why the smooth line is the
honest one --- rather than being told in a caption to trust it.
\end{column}
\end{columns}
\takeaway{The figure carries the argument the honesty box currently has to make in prose.}
\end{frame}

% =========================================================================
\begin{frame}{The basic picture is too bare}
\small
Basic currently gets a stripped-down still showing \textbf{only the radius}. Intermediate and
advanced get the full marked-up frame --- speed along the path, turn rate, and inward pull, each
drawn as its own arrow with its own label.
\vspace{0.6em}

Basic gets the fuller picture too, with every label kept in \textbf{plain words} rather than
symbols --- which is the rule that separates the basic level from the rest, and is untouched by
this change.
\vspace{0.6em}

\begin{tcolorbox}[colback=cPed!6,colframe=cPed,boxsep=2pt,left=6pt,right=6pt,top=3pt,bottom=3pt]
\small The picture a beginner sees should carry \emph{more} of the explanation than the one an
advanced reader sees, not less --- the beginner has the least text to fall back on.
\end{tcolorbox}
\takeaway{The figure is where a basic reader does most of the work, so it is the last thing that
should be stripped down.}
\end{frame}

% =========================================================================
\begin{frame}{Order of work, and what needs deciding}
\scriptsize
\renewcommand{\arraystretch}{1.45}
\begin{tabular}{@{}clp{7.4cm}@{}}
\toprule
& \textbf{work} & \textbf{why this order} \\
\midrule
\textbf{1} & the direction-sign error &
a correctness bug, in front of students now; needs no teaching decision \\
\textbf{2} & reading level of advanced; graph before equation &
named directly in the review; contained edits with an existing check to keep them honest \\
\textbf{3} & the two figure changes &
same, and independent of the restructure \\
\textbf{4} & sub-levels and bridges &
reshapes the tier architecture --- worth landing on material that is already correct and readable \\
\bottomrule
\end{tabular}
\vspace{7pt}

\textbf{\color{cAccent}Three open questions --- a teaching call, not a code fix:}
\begin{itemize}\setlength\itemsep{2pt}
  \item \textbf{Nine documents, or three documents with three graded sections inside each?}
        This decides how much of the pipeline moves.
  \item \textbf{What are the third steps of intermediate and advanced?} The proposal names the
        first two of each and stops.
  \item \textbf{``Show measurement and correction''} --- on the one clip whose camera angle is
        corrected, or on every graph that is smoothed?
\end{itemize}
\vspace{2pt}

Items 1--3 can be built and checked without the language model running, by replaying the
fixed part of an existing clip in about five minutes per run.
\takeaway{Fix what is wrong, lower the reading level, then restructure.}
\end{frame}

\end{document}
